% --- INICIO DE PLANTILLA PERSONALIZADA ---
\documentclass[oneside,11pt]{Latex/Classes/PhDthesisPSnPDF}

% --- Paquetes Originales ---
\usepackage{blindtext}
\usepackage{actuarialangle} 
\usepackage{tabularx}
\usepackage{float}
\usepackage{multirow, array}
\usepackage[usenames,dvipsnames,svgnames,table]{xcolor}
\usepackage{longtable}
\usepackage{latexsym,amsmath,amssymb,amsfonts}
\usepackage{enumitem}
\usepackage{xparse}
\usepackage{booktabs}
\usepackage[round, sort, numbers]{natbib}  
\usepackage{adjustbox}
\usepackage[utf8]{inputenc}
\usepackage{caption}
\usepackage{graphicx}
\usepackage{hyperref}
\usepackage{cleveref}

% Definición de colores
\definecolor{miblue}{RGB}{68, 114, 196}

% --- CAMBIO CRÍTICO 1: Usar \input para macros en el preámbulo ---
% \include da problemas antes del begin document. \input es seguro.
\input{Latex/Macros/MacroFile1}

% --- Definiciones de seguridad (por si la clase falla al cargar alguna variable) ---
\providecommand{\facultad}[1]{\def\lafacultad{#1}}
\providecommand{\escudofacultad}[1]{\def\elescudo{#1}}
\providecommand{\degree}[1]{\def\elgrado{#1}}
\providecommand{\director}[1]{\def\eldirector{#1}}
\providecommand{\lugar}[1]{\def\ellugar{#1}}
\providecommand{\degreedate}[1]{\def\lafecha{#1}}

% --- Configuración para bloques de código de R (Pandoc) ---

%%%%%%%%%%%%%%%%%%%%%%%%%%%%%%%%%%%%%%%%%%%%%%%%%%%%%%%%%%%%%%%%%%%%%%%%%%%%%%%%
%                         DATOS (Conectados a RMarkdown)                       %
%%%%%%%%%%%%%%%%%%%%%%%%%%%%%%%%%%%%%%%%%%%%%%%%%%%%%%%%%%%%%%%%%%%%%%%%%%%%%%%%
\title{Análisis del Catálogo de Amazon y Estrategia de Precios}
\author{Mariangel Morales \\ Camila Santiago}

% Lógica para Facultad: Si la defines en el YAML la usa, si no, usa la default

\facultad{Facultad de Ciencias Económicas y Sociales \\ Escuela de
Estadistica y Ciencias Actuariales}
  \escudofacultad{Latex/Classes/Escudos/faces}

\degree{Computación I}
\director{Jesus Ochoa \\ Oliver Triveño}
\degreedate{Abril 16 del 2026} 
\lugar{Caracas}

\portadatrue 

% Metadatos PDF
\keywords{}
\subject{}

% --- CORRECCIÓN PARA PANDOC NUEVO (Imágenes) ---
\makeatletter
\providecommand{\pandocbounded}[1]{#1}
\makeatother


%%%%%%%%%%%%%%%%%%%%%%%%%%%%%%%%%%%%%%%%%%%%%%%%%%%%%
%                   DOCUMENTO                       %
%%%%%%%%%%%%%%%%%%%%%%%%%%%%%%%%%%%%%%%%%%%%%%%%%%%%%
\begin{document}

\maketitle

%%%%%%%%%%%%%%%%%%%%%%%%%%%%%%%%%%%%%%%%%%%%%%%%%%%%%
%                  PRÓLOGO                          %
%%%%%%%%%%%%%%%%%%%%%%%%%%%%%%%%%%%%%%%%%%%%%%%%%%%%%
\frontmatter

% Puedes agregar dedicatorias desde el YAML si quieres, o dejarlas fijas aquí

%%%%%%%%%%%%%%%%%%%%%%%%%%%%%%%%%%%%%%%%%%%%%%%%%%%%%
%                   ÍNDICES                         %
%%%%%%%%%%%%%%%%%%%%%%%%%%%%%%%%%%%%%%%%%%%%%%%%%%%%%
\setcounter{secnumdepth}{3}
\setcounter{tocdepth}{3}

\tableofcontents



\cleardoublepage

  \cleardoublepage       % Empieza en página derecha (estándar en tesis)
  \chapter*{Resumen}     % Crea el título "Resumen" sin numeración (Capítulo *)
  \addcontentsline{toc}{chapter}{Resumen} % (Opcional) Lo añade al índice
  
  En el presente informe se realiza un análisis estadístico descriptivo
de las ventas de Amazon India, centrándose en su estructura del catalogo
y su estrategia de precios. Como principales hallazgos se indica que el
catálogo de Amazon Presenta una alta heterogeneidad con pocas categorías
dominantes qué concentran la mayoría de los productos, se identificó la
ausencia de correlación entre el precio original y el porcentaje de
descuento, lo que sugiere que no hay ninguna estrategia al momento de
aplicarle descuentos a los precio. Además, la detectaron valores
atípicos en los precios reflejan la diversidad económica del mercado
indio.             % Aquí se pega el texto que escribiste en el YAML
  
  \cleardoublepage

%%%%%%%%%%%%%%%%%%%%%%%%%%%%%%%%%%%%%%%%%%%%%%%%%%%%%
%                   CONTENIDO (El Corazón)          %
%%%%%%%%%%%%%%%%%%%%%%%%%%%%%%%%%%%%%%%%%%%%%%%%%%%%%
\mainmatter
\def\baselinestretch{1.5}

% --- CAMBIO CRÍTICO 2: Aquí RMarkdown inyecta todo tu texto ---
\chapter*{Introducción}

Desde la llegada de Amazon a India en junio de 2013. la empresa se ha
ido consolidando rápidamente en el comercio electrónico del país,
ofreciendo una amplia gama de productos y servicios a los consumidores
locales. Este trabajo tiene como propósito analizar la estructura del
catálogo Amazon a partir de un conjunto de datos de 1.465 productos, con
el objetivo de responder las interrogantes planteadas sobre la
concentración de categorías, la variabilidad de los precios y una
posible estrategia de descuento.

\chapter{Planteamiento del Problema}

Amazon es una multinacional tecnológica estadounidense, fundada por Jeff
Bezos en 1994, posteriormente, en 1995 se convierte en uno de los
primeros grandes minoristas en línea cuando empieza a vender libros de
forma online, y en 1996 se expande a la música, el vídeo y otros
productos, lanzando su primer sitio web no inglés en Canadá. Durante los
2000 se expande a nuevas categorías de productos, como juguetes y
electrónica, y lanza su sitio web en otros idiomas. Amazon sigue
creciendo y lanza su programa de afiliación Prime, que ofrece envíos
gratuitos ilimitados en dos días para millones de productos, además de
otras ventajas exclusivas (2003), así como también, lanza Kindle
e-lector y abrir una tienda física en Seattle (2006); su expansión
continúa con servicios en la nube, más variedad de productos, sucursales
en todo el mundo y Amazon music (2011). En el contexto de India, opera
bajo el nombre de Amazon India, donde los precios se expresan en rupias.

A pesar de su éxito comercial, su amplitud y variedad en el mercado
presenta dificultades para los vendedores, entre ellas, se señala que
Amazon tiende a priorizar sus propios productos frente a los de
terceros, además del elevado número de negocios que ofrecen otros
productos muy similares en la plataforma. Tanto los vendedores
independientes como sus competidores comercializan en el mismo
escaparate, sin margen de conveniencia para ninguna de las partes. La
principal virtud de la era digital son los datos. Para las empresas esta
es una valiosa información, pero Amazon no comparte estos datos con los
vendedores que distribuyen a través de su plataforma. Los vendedores se
encargan de llevar a cabo la transacción de los pedidos, pero no pueden
acceder a los clientes individualmente. Esto dificulta cualquier intento
de fidelización con sus consumidores; así, Amazon posee el monopolio de
los datos, y de las campañas promocionales. Un ejemplo desfavorable de
esto para los vendedores es la semana del Black Friday, durante las
ofertas vigentes, los distribuidores de forma obligatoria tienen que
ajustar sus ofertas y sus precios a lo establecido por la plataforma. En
este marketplace se paga un porcentaje por cada venta que se realice a
través de la plataforma, las comisiones de Amazon se encuentran entre el
7\% y el 15\% de la venta dependiendo del sector al que pertenezca el
artículo. Además, deben pagar al mes o por venta el plan elegido. Por
esta razón las marcas o negocios que venden productos que tienen un bajo
precio necesitan ampliar su catálogo o vender productos de mayor valor
para que les sea rentable vender en el marketplace.\\

En este sentido, se plantean las siguientes interrogantes:\\

\textbf{1.} ¿Existen categorías donde los descuentos son más altos que
el resto del catálogo y, si éstas existen, hay alguna relación entre el
precio original y su porcentaje de descuento?\\

\textbf{2.} ¿La distribución de los productos está equilibrada entre
categorías o, el catálogo es predominante y más voluminoso en otras
categorías?\\

\textbf{3.} ¿Dentro de las categorías, conviven productos de distintos
rangos de precios o existe una homogeneidad entre ellos?\\

\section{Justificación}

La presente investigación se justifica por su relevancia práctica y
académica. Por la parte práctica, los resultados les van a permitir a
los compradores, reconocer ofertas reales y engañosas, y a los
vendedores de amazon identificar patrones de descuento y así poder
ajustar sus estrategias de precios; y por la parte, de forma académica
se aporta una metodología replicable para el análisis de datos de
comercio electrónico utilizando python para el dashboard y R para el
informe técnico.\\

\section{Objetivos}

\subsection{Objetivo General}
\begin{itemize}
  \item{Describir la estructura del inventario de Amazon y analizar su estrategia de precios según sus diferentes categorías e identificar las brechas de precios y descuentos entre productos similares.}
\end{itemize}

\subsection{Objetivos Específicos}
\begin{itemize}
  \item{Analizar las diferencias entre productos similares, y la relación de precios entre el precio original y su porcentaje descuento en categorías con precios más altos.}
  \item{Identificar el número de productos existentes en cada categoría para determinar si el catálogo está distribuido equilibradamente o está concentrado en pocas categorías dominantes.}
  \item{Evaluar  la variabilidad de precios dentro de cada categoría para determinar si existen diferencias significativas entre productos de distintos rangos o si hay homogeneidad.}
\end{itemize}

\section{Variables}

Se tienen variables cualitativas nominales (id del producto, nombre del
producto, categoria) y cuantitativas continuas de razón (precio
original, porcentaje de descuento, precio con descuento).

\chapter{Marco Teórico}

\section{Bases Teóricas}

Las bases teóricas que se usarán en este trabajo para sustentar el
análisis e investigación de este tema son vitales para el entendimiento
de la misma, ya que, estas palabras por más ``comunes'' que puedan
sonar, sino se usan en el contexto adecuado pueden dificultar el
razonamiento del lector, las siguientes definiciones son:\\

\subsection{Catálogo}

Los catálogos son contenedores de elementos que están relacionados entre
sí a través de un contexto de negocio\\

\subsection{Rupias}

Es la moneda oficial de la India ``\,``, también está vinculada a otras
monedas, como el ngultrum butanés y la rupia nepalí.\\

\subsection{Estrategia de precios }

Enfoque estructurado que las empresas utilizan para determinar cuánto
cobrar por sus productos o servicios, se consideran diversos factores,
como la demanda del mercado, la competencia, los costos de producción y
la percepción de valor por parte del cliente.\\

\subsection{Marketplace }

Entorno virtual que facilita la compra y venta de productos o servicios
entre múltiples vendedores y compradores. Este tipo de plataforma actúa
como intermediario, reuniendo a vendedores y consumidores en un solo
lugar y ofreciendo una amplia variedad de productos o servicios en una
experiencia de compra unificada.\\

\subsection{Variabilidad de precios }

Esta variabilidad ayuda a las empresas a gestionar sus ingresos,
identifican qué productos contribuyen principalmente a sus beneficios
totales y cuáles requieren un enfoque de precios diferente para aumentar
los ingresos.\\

\subsection{Homogeneidad y heterogeneidad }

La homogeneidad se refiere a elementos semejantes, idénticos e iguales.
Por otro lado, la heterogeneidad se refiere a elementos diferentes,
distintos y no equivalentes.\\

\subsection{Porcentaje de descuento }

Número del descuento que posee el producto.\\

\subsection{Varianza  }

Medida de dispersión que representa la variabilidad de una serie de
datos con respecto a su media.\\

\subsection{Correlación  }

Es la asociación o relación entre dos variables numéricas que evalúa la
tendencia (aumento o disminución) en los datos.\\

\subsection{Vendedor }

ofrecen productos terminados o servicios a los compradores finales.\\

\subsection{Proveedor }

Suministran materias primas o componentes utilizados en la producción.\\

\subsection{Comprador }

Es la persona que realiza la transacción de compra, el término comprador
suele tener un matiz transaccional, ya que se refiere únicamente a la
acción de comprar.\\

\chapter{Marco Metodólogico}

En este capítulo se presentan los fundamentos estadísticos empleados que
hicieron posible analizar la base de datos, la cual fue obtenida de
Kaggle y se titula ``Amazon Sales'', contiene una parte del catálogo de
Amazon India; la misma, se conforma de 16 columnas y 1465 filas, se le
realizó una limpieza de columnas, dejando así solo 6 que serían las
variables clave para este análisis. Esta investigación adopta un enfoque
cuantitativo, de tipo descriptivo-correlacional.

\section{Metodología utilizada en el trabajo de investigación}

A través de la estadística descriptiva, se examinarán los datos para
responder las preguntas de investigación planteadas. Se utilizará la
media aritmética en este trabajo para saber el precio promedio por
categoría para el análisis de la estructura de precios del catálogo,
está se va a comparar con la mediana, para reconocer si hay productos
que afecten significativamente al promedio; también, se usarán medidas
como, el rango (podemos saber que tan amplia es la diferencia de
precios), rango intercuartílico (mide la dispersión del 50\% central de
los datos y se utiliza para identificar valores atípicos), desviación
típica (indica la dispersión de los datos, si es mayor, significa que
están más dispersos los datos), varianza (permite comparar la
variabilidad entre categorías con precios muy diferentes) y correlación
(si dos variables se relacionan de forma lineal o inversa).\\

Se van a cargar las librerias tidyverse, knitr, kableExtra, scales,
janitor y ggplot2 para el manejo de los datos, se carga el dataset para
la limpieza, Se eliminarán las columnas que no aportan al analisis de la
investigación, Se identifican los valores nulos del dataset, Se calculan
los cuartiles y el rango intercuartilico para identificar los datos
atipicos

\chapter{Análisis de Resultados}

\small Este capítulo presenta los resultados obtenidos al aplicar las
metodologías descritas en el capítulo anterior, se exponen las
respuestas a las preguntas planteadas y los resultados de los objetos
específicos, por medio de visualizaciones (creación de gráficos) que
ayudan a mostrar las distribuciones y relaciones de una forma mas
interactiva.

\section {¿Existen categorías donde los descuentos son más altos que el resto del catálogo y, si éstas existen, hay alguna relación entre el precio original y su porcentaje de descuento?}

\begin{center}\includegraphics[width=0.9\linewidth,]{informe_files/figure-latex/unnamed-chunk-3-1} \end{center}

\begin{center}\includegraphics[width=0.9\linewidth,]{informe_files/figure-latex/unnamed-chunk-3-2} \end{center}

\begin{center}\includegraphics[width=0.9\linewidth,]{informe_files/figure-latex/unnamed-chunk-3-3} \end{center}

Sí, hay categorías que notablemente poseen descuentos más altos, la
correlación que se distingue es negativa (inversa) y muy baja.\\

\section{¿La distribución de los productos está equilibrada entre categorías o, el catálogo es predominante y más voluminoso en otras categorías?}

\begin{center}\includegraphics[width=0.9\linewidth,]{informe_files/figure-latex/unnamed-chunk-4-1} \end{center}

No, el catálogo no esta esta equilibrado, notablemente por el gráfico se
concluye que su categoría mas predominante es ``Electronics''.\\

\section{¿Dentro de las categorías, conviven productos de distintos rangos de precios o existe una homogeneidad entre ellos?}

\begin{longtable}[t]{lrrr}
\caption{\label{tab:unnamed-chunk-5}Resumen de Precios por Categoría}\\
\toprule
Categoría & Precio Mínimo & Precio Máximo & Brecha\\
\midrule
\cellcolor{gray!10}{Car\&Motorbike} & \cellcolor{gray!10}{4000} & \cellcolor{gray!10}{4000} & \cellcolor{gray!10}{0}\\
Computers\&Accessories & 39 & 59890 & 59851\\
\cellcolor{gray!10}{Electronics} & \cellcolor{gray!10}{171} & \cellcolor{gray!10}{139900} & \cellcolor{gray!10}{139729}\\
Health\&PersonalCare & 1900 & 1900 & 0\\
\cellcolor{gray!10}{Home\&Kitchen} & \cellcolor{gray!10}{79} & \cellcolor{gray!10}{75990} & \cellcolor{gray!10}{75911}\\
\addlinespace
HomeImprovement & 599 & 999 & 400\\
\cellcolor{gray!10}{MusicalInstruments} & \cellcolor{gray!10}{699} & \cellcolor{gray!10}{1995} & \cellcolor{gray!10}{1296}\\
OfficeProducts & 50 & 2999 & 2949\\
\cellcolor{gray!10}{Toys\&Games} & \cellcolor{gray!10}{150} & \cellcolor{gray!10}{150} & \cellcolor{gray!10}{0}\\
\bottomrule
\end{longtable}

Sí, existen categorías donde la diferencia de precios es alta, se ve
heterogeneidad\\

\chapter{Conclusiones y Recomendaciones}

\textbf{1.-} Conclusion 1 del trabajo de investigación.\\

\textbf{Recomendación} Recomendación 1 del trabajo de investigación.\\

\textbf{2.-} Conclusion 2 del trabajo de investigación.\\

\textbf{Recomendación} Recomendación 2 del trabajo de investigación.\\

\textbf{n.-} Conclusion n del trabajo de investigación.\\

\textbf{Recomendación} Recomendación n del trabajo de investigación.\\

\appendix
\renewcommand{\thechapter}{\Alph{chapter}}

\chapter{Bibliografias}

Prysync:
\url{https://prisync.com/es/diccionario-de-precios/variacion-de-precios/}\\

Statistic:
\url{https://www.statisticshowto.com/homogeneity-and-heterogeneity-in-statistics/}\\

Questionpro: \url{https://www.questionpro.com/blog/es/varianza/}\\

Inbestme:
\url{https://www.inbestme.com/es/es/blog/que-es-la-correlacion//}

Translate Google:
\url{https://ramp-com.translate.goog/blog/vendor-vs-supplier?_x_tr_sl=en&_x_tr_tl=es&_x_tr_hl=es&_x_tr_pto=tc&_x_tr_hist=true}\\

Blog empresas:
\url{https://blogempresas.yoigo.com/usuario-consumidor-cliente-comprador-son-lo-mismo//}

Arimetrics: \url{https://www.arimetrics.com/glosario-digital/amazon}\\

Cyberclick:
\url{https://www.cyberclick.es/numerical-blog/ventajas-y-desventajas-de-vender-en-amazon-marketplace}\\

Globalshopaholics:
\url{https://globalshopaholics.com/es/blog/pros-y-contras-de-comprar-en-amazon/}\\

Epinium:
\url{https://epinium.com/es/blog/estrategiapreciosventa-amazon/\#}:\textasciitilde:text=lanzar\%20el\%20producto.-,\%C2\%BFCu\%C3\%A1l\%20es\%20la\%20estrategia\%20de\%20precios\%20de\%20venta\%20en\%20Amazon,ganar\%20la\%20\%E2\%80\%9CBuy\%20Box\%E2\%80\%9D.\\

%%%%%%%%%%%%%%%%%%%%%%%%%%%%%%%%%%%%%%%%%%%%%%%%%%%%%
%                   REFERENCIAS                     %
%%%%%%%%%%%%%%%%%%%%%%%%%%%%%%%%%%%%%%%%%%%%%%%%%%%%%
% Mantenemos la estructura natbib de tu plantilla
\renewcommand{\bibname}{Lista de Referencias}
\bibliographystyle{apalike} 
\bibliography{bibliografia.bib} 

%%%%%%%%%%%%%%%%%%%%%%%%%%%%%%%%%%%%%%%%%%%%%%%%%%%%%
%                   APÉNDICES                       %
%%%%%%%%%%%%%%%%%%%%%%%%%%%%%%%%%%%%%%%%%%%%%%%%%%%%%

\end{document}
